\cotefile{93}

\scanpage{193}

que lorsque il n'est pas nul, il \uline{n'est pas identique} au groupe des $\omega$ dont un multiple est logarithmique, i.e. il y a des \deleted{\unsure{des str.}} connexions intégrables sur $\mathcal O_X$ qui ne sont pas isomorphes à la connexion standard, [en fait le groupe des classes d'isomorphie de telles connexions intégrables est isomorphe à \deleted{$\mathbb Q/\mathbb Z$} $L\otimes\mathbb Q/\mathbb Z\simeq(\mathbb Q/\mathbb Z)^{r}$]. Par contre, si $X$ est complet, \inserted{donc $\Gamma(X,\mathcal O_X^{*})$ \deleted{\unsure{formé seul.}} formé de sections horizontales,} \unsure{tte|une} connexion intégrable sur $\mathcal O_X$ est triviale, i.e. une $1$-forme différentielle qui est loc. ét une dérivée logarithmique, est nulle. \deleted{La classification des modules inversibles avec \uline{connexion à courbure} nulle, i.e. à stratification, est donc}

Revenons au cas où $k$ est de car. quelconque. Je dis que si $\Gamma(X,\mathcal O_X)^{*}$ est formé de fonctions horizontales, \deleted{\unsure{des}} p.ex. $X$ propre sur $k$ et séparable sur $k$, alors sur $\underline O_X$ il n'y a comme seule \deleted{\unsure{structure}} stratification \uline{intégrable} que la structure habituelle. En effet, $\underline O_X^{\otimes n}=\underline O_X$ \inserted{($n$ premier à $p$)} a une section horizontale \inserted{$\sigma$, pr $\unsure{n\ne0|n\gg0|\text{tout}\ne0}$}, \deleted{\unsure{donc qui est}} pour la stratification $T^{\otimes n}$, \deleted{qui est aussi horizontale}. Donc la section $\sigma_0=1$ de $\underline O_X$ satisfait $\sigma_0^{\otimes n}=\lambda\sigma$, $\lambda\in\Gamma(X,\underline O_X)^{*}$, \deleted{donc} et $\lambda$ horizontale, \inserted{par hyp.,} donc $\lambda\sigma$ aussi, donc $\sigma_0^{\otimes n}$ horizontale, donc $\sigma_0$ aussi. // Par suite, on voit que la classification des \deleted{\unsure{fibr.}} modules inversibles à connexion intégrable est, dans ce cas, \uline{identique} à celle des modules inversibles \deleted{\unsure{dont}} $L$ tels que $\exists\ n>0$, $n$ premier à $p$, avec $L^{\otimes n}\simeq\underline O_X$

% archive sheet number: 192
\scanpage{194}

\uline{Questions}

$S$ schéma de type fini sur $\mathbb Z$, dominant $\mathbb Z$, $S$ intègre pt gén. $\eta$

$X$ schéma lisse de type fini sur $S$.

$E$ module localmt libre sur $X$, avec connexion à courbure nulle relativement à $S$.

Le \inserted{module à} connexion \inserted{$E_{\eta}$ sur $X_{\eta}$} est dit \uline{de niveau zéro} \inserted{ou \uline{intégrable}} \deleted{si l'on peut trouver, sur $X_{\eta}$, un} l'on peut trouver un revêtement fini étale surjectif de $X_{\eta}$ sur lequel $E_{\eta}$ devient localmt trivial, i.e. ``intégrable'' au sens (Zar). \deleted{Une condition}\inserted{\unsure{Il revient|Je veux}} \deleted{nécessaire pour ceci est que} au \unsure{$\simeq$|$=$} de dire que sur une extension alg. close \deleted{\unsure{finie}} $K$ de $k(\eta)$, $E_K$ sur $X_K$ soit associé à des repr. linéaires des groupes fondamentaux des composantes connexes de $E_K$. Et si on choisit $K=\mathbb C$, que les représentations des $\pi_1$ \inserted{transcendants} des composantes \deleted{\unsure{gén.}} connexes, définies par $E_{\mathbb C}$, \deleted{\unsure{soit}} se fassent à travers des groupes finis.

% slanted

\margin{définition a un sens sur n'importe quelle \deleted{\unsure{schéma}} base --- Il faut mieux de dire : localmt \deleted{\unsure{simplemt}} trivial (pour la top. étale) ...}

Cette notion peut se varier aussi pour n'importe quel schéma en groupe de \unsure{t.f.|déf.|dif.|type f.} $G$ sur $S$, on prend sur $X$ un torseur $P$ sous $G_X$, avec connexion à courbure nulle rel/$S$.

La question qui suit pourra se poser aussi pour des tels groupes.

Remarquons que si la connexion de $E$ est intégrable, alors

% archive sheet number: 193
\scanpage{195}

il existe un \inserted{voisinage} ouvert \deleted{\unsure{de}} $U$ \deleted{$\ne\varnothing$ de $S$} de $\eta$ dans $S$, tel que pour tt $s\in U$, \inserted{$k(s)$ de car. $p>0$,} \deleted{$E_s$ sur $X_s$, avec la connexion relativement à $k(s)$} $p$-courbure de $E_s$ sur $X_s/k(s)$ soit nulle. \inserted{ou encore, intégrables ---} (Cette dernière condition est d'ailleurs équivalente à la condition analogue, \deleted{où} en se limitant aux pts \uline{fermés} de $U$, mais peu importe].

\uline{Question} La \uuline{réciproque est-elle vraie} ?

\uline{Cas \deleted{\unsure{des}} où} $E$ \uline{est de rang} $1$. Notons que \deleted{\unsure{si}} $E_{\eta}$ intégrable signifie que $E$ est associé à un revêtement \unsure{Kummerien|Kummérien} de $X_{\eta}$, i.e. \deleted{obtenu en un module $L$} sa connexion est associé à une trivialisation de $E^{\otimes n}$, pour $n\geq1$ convenable.

La question précédente se décompose en deux :

% vertical rule at left of a)

a) Si la connexion des $E_s$ \inserted{($s\in S$, $s$ de car $>0$)} est intégrable, est-il vrai que l'image de $E$ dans $\mathrm{Pic}(X_\eta)$ est de torsion ?

\deleted{b)} Si la réponse est affirmative, on est ramené dans le Pb au cas où $E\simeq\underline O_X$ (comme ne uniquement, pas la connexion, à priori).

% vertical rule at left of b)

b) Si $\omega$ est une $1$-forme différentielle \inserted{[fermée]} sur $X/S$, et si les $\omega_s$ sur les $X_s/s$ ($s\in S$, $s$ de car $>0$) sont \uline{loc.} des différentielles logarithmiques, alors est-il vrai que \inserted{(un multiple entier de)} $\omega_\eta$ est une diff. \uline{log}arithmique ?

\margin{ce qui se vérifie par l'opération de Cartier $\widetilde C\omega_s=\widetilde\omega_s$}

% archive sheet number: 194
\scanpage{196}

(resp. de rang unipotent nul, resp. affine, resp. unipotent).

\uline{NB} Les cas ``\inserted{connexe} affine'', ``unipotent'' sont triviaux, et essentiellement triviaux \deleted{(le deuxième, \unsure{tout|tous|base} au moins, si l'on \unsure{connait} plus \unsure{générlt|gén.}} (car le premier des deux est \unsure{obvie|évident|achevé|autre}), grâce au th. de Chevalley. Noter que dans le cas unipotent, l'hypothèse signifie que $\omega(s)$ est d'ordre $p^{\unsure{\infty|n}}$ ($p$ la car. en $s$). Le cas \deleted{\unsure{par}} ``rang unipotent nul'' peut se reformuler ainsi, lorsque $E$ est commutatif :

% vertical rule beside the following assertion

Supposons qu'il existe un entier $N$ tel que, pour tt point \uline{fermé} $s$ de $S$, $\omega(s)$ soit \deleted{\unsure{divis.}} dans

\[
(p^NE)(s)=(p^NE_s)(s)=p^NE(s),
\]

i.e. qu'il soit divisible par $p^N$. Alors $\omega(\eta)$ appartient à un sous-groupe algébrique de $E_{\eta}$ de rang \deleted{\unsure{uni}} unipotent nul, donc une puissance de $\omega(\eta)$ appartient à un ss-gpe extension d'une VA par un tore.

% slanted

\margin{Pour ce pb, on peut supposer que $E_\eta$ n'a pas de sous-gpes connexes de rang unipotent nul sauf zéro, \deleted{\unsure{de qq façon|de tte façon}} i.e. \unsure{prend|plus} une extension d'une V.A. par un vectoriel, qui ne splitte sur aucune sous variété abélienne $\ne0$, et il faut conclure que $\omega$, telle que ci-contre, est nécessairement \uline{d'ordre fini}.}

\deleted{Si} ce résultat étant acquis, \unsure{le|ce|les} \deleted{\unsure{question}} cas ``connexe propre'' se ramènerait au cas où $E$ est déjà une extension d'un schéma abélien $A$ par un tore $T$.

\deleted{Si p.ex. $E$ est une extension \uline{non triviale} de \deleted{\unsure{groupe}}\inserted{L'image} \unsure{la} \deleted{\unsure{cas d'une}} \inserted{infini d'} une variété abélienne $A$ par un tore, sur disons, alors il faudrait voir que si $\omega$ est telle que $\omega(s)$} Contrairement aux trois autres cas, ce cas-ci semble également intéressant si on suppose $S$ de car. $p>0$ \unsure{lui-$=$|lui-même} !

% archive sheet number: 195
\scanpage{197}

\deleted{\unsure{Abordons|Considérons} en particulier le cas où $X\to S$ est propre. \unsure{Donc on a|Dans ce cas} $\underline{\unsure{W|\omega}}=f_*(\Omega^1_{X/S})$ faisceau coh. sur $S$, et quitte à restreindre $S$, on peut supposer que}

\[
\deleted{\underline{\underline{W}}=f_*(Z^1_{X/S})=\operatorname{Ker}\bigl(f_*(\Omega_{X/S})\longrightarrow f_*(\Omega^2_{X/S})\bigr).}
\]

Dans le cas où $X\to S$ est un schéma abélien, introduisant le schéma abélien dual $A$, et son extension vectorielle universelle $E$, la question \deleted{\unsure{se|est}} équivaut à la suivante : \deleted{Soit $\varphi$ une section}

\margin{cette traduction vaut, plus gén., si $X\to S$ propre et lisse}

% vertical rule beside c)

c) Soit $\deleted{\varphi}\omega$ une section de $E$ sur $S$. Supposons que pour tout pt \uline{fermé} $s\in S$, $\omega(s)\in A_s^{\unsure{(p/k(s))|(p/k_s)|(p(k(s)))}}$ \inserted{$(=pE_s)$}, où $p=\mathrm{car}\,k(s)$. Alors la \deleted{\unsure{l'image de|image de}}\inserted{section} $\omega$ de $\Gamma(S,A_s)$ est-elle d'\uline{ordre fini} ? [NB. a) signifie que l'image de $\omega$ dans $\Gamma(S,A_S)$ est d'ordre fini, b) signifie que si on \inserted{suppose de plus que} $\omega$ est une section de $V=\mathrm{Ker}(E\to A)$, alors $\omega$ est nulle].

Ceci suggère la question

d) Soit $S$ de type fini sur $\mathbb Z$, intègre, \inserted{de pt gén $\eta$}, dominant $\mathrm{Spec}\,\mathbb Z$, \deleted{\unsure{Soit}} $E$ un schéma en groupes \deleted{\unsure{de loc}} de t.f. sur $S$ (lisse connexe si on veut\unsure{ !|}), $\omega$ une section de $E$. Supposons que pour tt pt fermé $s\in S$, $\omega(s)$ appartienne \deleted{au ss-gpe} à un sous-gpe de $E_s$ qui est \uline{propre} \inserted{\uline{connexe}} sur $k(s)$ [resp. \inserted{connexe} de rang unipotent nul, resp. \inserted{\uline{connexe}} affine, resp. \inserted{connexe} unipotent] alors $\omega(\eta)$ appartient à un ss-gpe de $E_\eta$ qui est \uline{propre}

% archive sheet number: 196
\scanpage{198}

La partie a) dans c) suggère la question suivante :

e) Soit $A$ schéma abélien sur $S$, $\omega$ section de $A$, telle que pour tt pt fermé $s\in S$, $\omega(s)$ se remonte à $A_s^{(p)}$ via Verschiebung $A_s^{(\unsure{p/s|p/k|p/S})}\longrightarrow A_s$. Est-il vrai qu'alors $\omega$ est d'ordre fini ?

\deleted{Dans le cas où $S$ lui-\unsure{$=$|$\simeq$} est de car $p>0$ (au lieu de dominer $\mathrm{Spec}\,\mathbb Z$) on dispose de $A^{p/S}\longrightarrow A$, et l'image inverse de $\omega$ dans $A^{p/S}$ est un revêtt \unsure{de|ét}}

\margin{fin 1\textsuperscript{ere} partie,}

% archive sheet number: 197
\scanpage{199}

Soit $X$ schéma sur $S$, tel que l'on soit sous les conditions permettant de définir $\widetilde C:f_{X/S*}(\Omega^1_{X/S})\longrightarrow\underline{\Omega}^1_{X^{p/S}/S}$, d'où

\margin{On suppose donc $f_*(\underline{\Omega}^1_{X/S})=f_*(Z^1_{X/S})$ et $\Lambda\simeq\underline O_{X^{p/S}}$}

\[
 f_*(\underline\Omega^1_{X/S})\overset{\mathrm{d\acute ef}}{=}\underline W\longrightarrow
 f_*(\underline\Omega^1_{X^{p/S}/S})\stackrel{\sim}{\longleftarrow}
 f_*(\Omega^1_{X/S})^{(p)}=\underline\omega^{(p)}.
\]

\margin{sous réserve que ...}

\margin{\uline{NB} au lieu de $f_*$, on pourrait prendre $\mathrm R^if_*$}

Donc on trouve un hom $\mathcal O_S$-linéaire

\[
\widetilde C:\quad\underline\omega\mathbin{\unsure{-|\longrightarrow}}\underline\omega^{(p)}
\]

on peut former alors $\widetilde C^{(p^i)}:\underline\omega^{(p^i)}\longrightarrow(\underline\omega^{(p)})^{(p^i)}=\underline\omega^{(p^{i+1})}$, et on trouve un composé $\mathcal O_S$-linéaire

\[
F_*^{(\nu)}=\underbrace{\widetilde C^{(p^{\nu-1})}\widetilde C^{(p^{\nu-2})}\ldots
\widetilde C^{p}\widetilde C}_{\nu\ \text{facteurs}}:
\quad\underline\omega\longrightarrow\underline\omega^{(p^\nu)}
\]

Si $\lambda\longrightarrow\lambda^{p^\nu}$ sur $\underline O_S^{*}$ est l'identité, p.ex si $S=\mathrm{Spec}\,\mathbb F_{p^\nu}$, on trouve donc que \deleted{\unsure{le}}

\[
F_*^{(\nu)}:\quad\underline\omega\longrightarrow\underline\omega
\]

est un \uline{endomorphisme} \inserted{$\mathcal O_S$-linéaire} de $\underline\omega$.

Soit $\omega$ une section de $\underline\omega$ telle que l'on ait

\[
\widetilde C\omega=\omega^{(p)}
\]

(ce qui correspond à une \deleted{section} $1$-forme $\omega$ sur $X$ telle que $\widetilde C\omega=\widetilde\omega$ i.e. qui est loc. une dérivée logarithmique ....). On conclut, par récurrence

\[
\deleted{\widetilde C^{\unsure{(p^i)|(p^\nu)}}}F^{(\unsure{i|\nu})}\omega=\omega^{(p^{\unsure{i|\nu}})},
\]

et dans le cas où $\lambda=\lambda^{p^\nu}$ sur $\underline O_S^{*}$, on conclut que l'on a

\[
\boxed{F_*^{(\nu)}\omega=\omega}
\]

% archive sheet number: 198
\scanpage{200}

\deleted{Soit $X$ schéma sur $S$, tel que $X\longrightarrow X^{p/S}$ un épimorphisme (p.ex. $X$ lisse sur $S$, plus gén. \deleted{il existe \unsure{localmt|loclt} une \unsure{$p$-base} relative à}\inserted{$X\longrightarrow S$ est séparable loc de prés finie,} et \deleted{$X/S$, et les fibres géométriques, les fibres géom. ont un} $\underline\Omega^1_{X/S}$ est loc. libre}

% archive sheet number: 199
\scanpage{201}

Notons que l'on peut aussi considérer

\[
f_{X/S}:X\mathbin{\unsure{-|\longrightarrow}}X^{p/S}
\]

d'où, en itérant, des

\[
f_{X/S}^{p^i/S}:X^{p^i/S}\longrightarrow X^{p^{i+1}/S}
\]

et par suite un composé

\[
F_{X/S}^{(\nu)}=f_{X/S}^{(p^{\nu-1})}f_{X/S}^{(p^{\nu-2})}\ldots f_{X/S}^{\deleted{\unsure{p}}}:
X\mathbin{-}X^{(p^\nu/S)}
\]

et on aura, si $X/S$ est \uline{lisse} \uline{de dim relative} $1$

\[
F_*^{(\nu)}=\operatorname{Tr}_{F_{X/S}^{(\nu)}}=\underset{\text{Gysin}}{\bigl(F_{X/S}^{\nu}\bigr)_*}
\]

en identifiant (sous réserve de formule de chgt. de base) $\deleted{\unsure{X^{p^i}}}$\deleted{\unsure{$f_*$}} $g_{\nu*}(\underline\Omega^1_{X^{p^\nu/S}})$ avec $g_{0*}(\Omega^1_{X/S})^{(p^{\unsure{\nu|1}})}=\underline\omega^{(p^\nu)}$. On trouve en particulier, si $\lambda\longmapsto\lambda^{p^\nu}$ est l'identité sur $\mathcal O_S^{*}$, donc $X^{p^\nu/S}\overset{\mathrm{can}}{\simeq}X$, que

\[
F^\nu_{X/S}:X\mathbin{-}X
\]

est un $S$-endomorphisme, qui d'ailleurs n'est autre que l'itéré $\nu$\textsuperscript{ième} de $F_X^\nu$ (compte tenu que par hyp. $F_S^\nu=\mathrm{id}_S$, \unsure{de sorte|du fait} que $F_X^\nu$ est bien un $S$-endomorphisme) [En général, et sans hyp. sur $S$, on peut dire que le composé

\diagram{\begin{tikzcd}[ampersand replacement=\&]
X\arrow[r,"F_{X/S}^{(\nu)}"] \& X^{(p^\nu/S)}\arrow[r,"\mathrm{can}"]\arrow[d,dotted] \& X\arrow[d,dashed]\\
\& {\bigl[\,S}\arrow[r,"{F_S^\nu}"] \arrow[ur,phantom,"\mathrm{cart}" description] \& {S\,\bigr]}
\end{tikzcd}}

est $F_X^\nu$ ]

% archive sheet number: 200
\scanpage{202}

\deleted{Supposons qu'on ait $S=\mathrm{Spec}\,\mathbb F_{p^\nu}$. Alors on peut définir $\deleted{\unsure{f}}F_X^{(\nu)}$, qui est un $S$-endomorphisme, et qui agit donc sur $\underline\omega$.}

% archive sheet number: 201
\scanpage{203}

Et on trouve donc que l'endomorphisme $F_*^{(\nu)}$ \deleted{de} $\omega$ n'est autre que l'\uline{hom trace} \inserted{$F_{X/S*}^{(\nu)}$ \deleted{\unsure{$=$}}}, associé à l'endom $F_{X/S}^{(\nu)}=(F_X)^\nu$ de $X/S$, [qui est transposé de l'endomorphisme $F_{X/S}^{(\nu)*}$ sur $\mathbb R^{n-i}f_*(\underline\Omega^{n-1}_{X/S})$, quand $X$ est propre et lisse sur $S$ --- ].

\deleted{Plaçons nous dans le cas} Notons que nous avons la suite exacte

\diagram{\begin{tikzcd}[ampersand replacement=\&,column sep=small]
0\arrow[r] \& \underline\omega\arrow[r]\arrow[d,equal] \& \mathcal H^1_{\mathrm{DR}}(X/S)\arrow[r]\&\deleted{\unsure{R}}\mathrm R^1 f_*(\underline O_{X/S})\\
\& \mathcal H^{1,0}_{\deleted{\unsure{X/S|\mathrm{DR}}}}(X/S) \& \&
\end{tikzcd}}

[avec un $0$ à droite dans les cas favorables), et que l'hom de Gysin \inserted{sur $\underline\omega$} associé à $F_{X/S}^{(\nu)}$ est compatible avec un hom de $=$ nom sur $\mathcal H^1_{\mathrm{DR}}(X/S)\longrightarrow\mathcal H^1_{\mathrm{DR}}(X^{p^\nu/S})$ \inserted{ss réserve $\unsure{\simeq|H}\ \mathcal H^1_{\mathrm{DR}}(X/S)^{(p^\nu)}$}, (qui \deleted{\unsure{sont|est}} des endom si $S=\mathrm{Spec}\,\mathbb F_{p^\nu}$). \deleted{Donc}

D'ailleurs, dans ce cas, \inserted{et si $X/S$ est propre,} \deleted{sauf erreur}, cet hom de Gysin induit zéro sur $\mathcal H^1_{\mathrm{DR}}(X/S)/\underline\omega\simeq\mathrm R^1 f_*(\underline O_{X/S})$, \deleted{il est peut-être \unsure{$=$|même|$\equiv$} nul sur les} (\unsure{regarder|raisonner|conjecture} par transposition), donc si $S=\mathrm{Spec}\,\unsure{\mathbb F_{p^\nu}|\#p^\nu}$, ses valeurs propres dans $\mathcal H^1_{\mathrm{DR}}(X/S)$ \deleted{et celles qui sont $\ne0$} sont aussi celles dans $\underline\omega$, et \deleted{\unsure{on}} \inserted{celles qui sont $\ne0$} sont aussi les valeurs propres \inserted{mod $p$} de Frobenius $\ell$-adique ($\ell\ne p$) \deleted{\unsure{sont cristallin|ou cristallin}}, ou cristallin qui sont des unités $p$-adiques. (Katz),

On voit donc que les $\omega$ tels que $\widetilde C\omega=\widetilde\omega$

% archive sheet number: 202
\scanpage{204}

i.e. qui sont loc. des dérivées logarithmiques, permettent de construire des \deleted{\unsure{valeurs|vecteurs}}\inserted{\unsure{vecteurs|valeurs}} propres de $F_{\deleted{X/S}}^{\nu}$ \deleted{qui se \unsure{x} de \unsure{$\omega$}} ($\ell$-adique, ou cristallin) qui sont $\equiv1$ $(p)$.

Plus précisément, il existe $\omega$ avec $\widetilde C\omega=\widetilde\omega$, $\omega\ne0$, ssi il existe une valeur propre de Frobenius, congrue à $1$ mod $p$. La dim de l'espace \deleted{des} engendré par les $\omega$ (\uuline{\unsure{vect.|vectr.} sur} $k$ !) tels que $\widetilde C\omega=\widetilde\omega$ est égale au nb des valeurs propres de Frobenius congrues à $1$ mod $p$. C'est aussi égal à la dim sur $\mathbb F_p$ du groupe \deleted{des \unsure{pts|pt} de la jacobienne} $ {}_pA(k)$, $A$ la jacobienne.

\margin{Fin 2\textsuperscript{e} partie ..}

% archive sheet number: 203
\scanpage{205}

$\left\{\begin{gathered} \text{\uline{Variantes connexifiantes}}\\ \text{\uline{des foncteurs de}}\\ \text{\uline{Picard relatifs}} \end{gathered}\right.$

\uline{\unsure{Katz|Notes|Att}. .}

% archive sheet number: 204
\scanpage{206}

Soit $f:X\longrightarrow S$ un morphisme de schémas, $f^{\mathrm{pl}}$ le morphisme correspondant des sites fppf.

\textcircled{1} $P^c=P^c_{X/S}:\ (\mathrm{Sch})_{/S}^{\unsure{o|\circ}}\longrightarrow(\mathrm{Ab})$, \deleted{défini par}\inserted{faisceau associé au préfaisceau} \deleted{$P^c(S')=$} $S'\longmapsto$ classes de modules inversibles sur $X_{S'}$, munis d'une connexion rel/$S'$

\[
0\longrightarrow\frac{f_*^{\mathrm{pl}}(\underline\Omega^1_{X/S})}{\mathrm{Im}\, f_*^{\mathrm{pl}}(\underline O^*_{X/S})}
\longrightarrow P^c_{X/S}\longrightarrow\underline{\mathrm{Pic}}_{X/S}
\longrightarrow\mathrm R^1 f_*^{\mathrm{pl}}(\underline\Omega^1_{X/S})\tag*{(1.1)}
\]

\margin{$f_*^{\mathrm{pl}}\bigl(\underline O^*_{X/S}\bigr)\longrightarrow f_*^{\mathrm{pl}}(\Omega^1_{X/S})$ donné par $\varphi\longmapsto d\varphi/\varphi$.}

% slanted

\margin{Cet hom est nul si $f$ est propre et séparable, \inserted{de prés. finie} p.ex. propre et lisse}

\uline{NB} Si $f$ \deleted{propre et \unsure{gén. sép.|plat}}\inserted{cohérent}, et $\underline\Omega^1_{X/S}$ cohom. plat/$S$ en dim $0$, ou $\mathrm R^1 f_*(\underline\Omega^1_{X/S})$ commute au chgt de base, alors pour tt schéma abélien $A$ sur $S$, tout hom $B\mathbin{\unsure{-|\longrightarrow}}\mathrm R^1 f_*^{\mathrm{pl}}(\underline\Omega^{1}_{X/S})$ est nul. donc si on a un hom $B\overset{u}{\longrightarrow}\underline{\mathrm{Pic}}_{X/S}$, on a $u(B)\subset\mathrm{Im}(P^c\longrightarrow\underline{\mathrm{Pic}}_{X/S})$, \unsure{d'où|donc} suite exacte

\[
\cdot\mathbin{-}f_*^{\mathrm{pl}}(\Omega^1_{X/S})\longrightarrow u^*(P^c)\longrightarrow B\longrightarrow1
\tag*{\unsure{(1.2')|(1.2)}}
\]

Ceci s'applique en particulier si \uline{$f$} propre sur $S$, \deleted{\unsure{do}} donc $\underline{\mathrm{Pic}}^0_{X/S}$ est défini, et si $\boxed{B=\underline{\mathrm{Pic}}^0_{X/S}}$ est un schéma abélien. Alors, $u$ étant l'inclusion, on aura $u^*(P^c)=(P^c)^0$, et on trouve suite exacte

\[
\cdot\longrightarrow f_*^{\mathrm{pl}}(\underline\Omega^1_{X/S})\longrightarrow(P^c_{X/S})^0
\longrightarrow\underset{\substack{\Vert\\\underline{\mathrm{Pic}}^0_{X/S}}}{B}\longrightarrow\cdot
\tag*{\unsure{(1.2)|(1.2')|(1.3)}}
\]

\uline{Re NB} \unsure{$\alpha_1$|ceci|ce}, et s'il est défini, \unsure{$\alpha_2$|ceci|ce|cas}, est fonctoriel en $X/S$, et commute au chgt de base.

\textcircled{2} $P^{ci}=P^{ci}_{X/S}$, sous-Groupe de $P^c_{X/S}$, faisceau associé au préfaisceau $S'\longmapsto$ classes de modules inversibles sur $X_{S'}$, munis d'une connexion rel. à $S'$ qui est \uline{à courbure nulle}

% archive sheet number: 205
